\section{複素数の積と行列}\label{節：複素数の積と行列}

複素数$u+iv$を$x+iy$に掛けると
\[
 (u+iv)(x+iy)=(ux-vy)+i(vx+uy).
\]
実部と虚部を縦に並べれば,
\[
 \begin{pmatrix}x'\\y'\end{pmatrix}
 =
 \begin{pmatrix}u&-v\\v&u\end{pmatrix}
 \begin{pmatrix}x\\y\end{pmatrix}.
\]
$u=\rho\cos\theta,\,v=\rho\sin\theta$なら,\,これは
\[
 \rho
 \begin{pmatrix}
 \cos\theta&-\sin\theta\\
 \sin\theta&\cos\theta
 \end{pmatrix}
\]
である. 複素数の掛け算と,\,二成分の実数の行列更新が一致した.

この行列の特性方程式は
\[
 X^2-2uX+(u^2+v^2)=0
\]
で,\,根は$u\pm iv$である.
実数の行列に非実数の固有値が現れることは,\,回転の倍率を複素数で記録したことに対応する.

ただし非実数の固有値を持つ一般の実行列が,\,元の座標でこの形に見えるとは限らない.
座標を変えて初めて回転と拡大・縮小になる場合もある.
その変換を次節で具体的に作る.


\[
 z_{n+1}=\lambda z_n,\qquad z_n=\lambda^{n-1}z_1
\]
は複素数の等比数列であり,実部と虚部を分ければ実数の連立漸化式である.
実数の行列と複素数は,この場合には同じ更新を二通りに書いている.
次節では,隣接3項間漸化式からこの形の複素数を作る.
